\documentclass{classeENS}
\usepackage[utf8]{inputenc}
\usepackage[T1]{fontenc}
\usepackage{geometry}
\usepackage{graphicx}
\usepackage{float}
\usepackage{hyperref}
\usepackage[english,indonesian]{babel}

\title{Optimasi Pemrosesan Template LaTeX Menggunakan Pendekatan Berbasis AI}
\author{Academic Support Agent}
\date{\today}

\begin{document}
\maketitle

\begin{abstract}

\end{abstract}

\chapter{PENDAHULUAN}


\chapter{TINJAUAN PUSTAKA}


% Original content (commented out):
% \documentclass{classeENS}
\usepackage{lipsum}

\title{Template rapport ENS } %Titre du fichier

\begin{document}

%----------- Informations du rapport ---------

\titre{Titre du rapport} %Titre du fichier .pdf
\departement{Département de Physique} 
\matiere{Matière} 
\motif{TP/Stage}

\enseignant{Prénom \textsc{Nom}} %Nom de l'enseignant
%\tuteurpcp{}
\eleves{Prénom \textsc{Nom}} %Nom des élèves

%----------- Initialisation -------------------
        
\fairemarges %Afficher les marges
\fairepagedegarde %Créer la page de garde
\tabledematieres %Créer la table de matières
\tabledefigure

%------------ Corps du rapport ----------------

\section{Introduction}

\section{Conclusion}

\end{document}

\end{document}